% Anthropic: Research Engineer, RL and post-training. Not served on the website.
\documentclass[10pt]{article}
\usepackage[margin=0.5in,vmargin=0.36in]{geometry}
\usepackage{fontspec}\setmainfont{Times New Roman}
\usepackage{enumitem}
\usepackage{titlesec}
\usepackage[hidelinks]{hyperref}
\hypersetup{pdftitle={Amal Mehta — Resume}, pdfauthor={Amal Mehta}}
\pagestyle{empty}
\setlength{\parindent}{0pt}
\titleformat{\section}{\large\bfseries}{}{0pt}{}[\vspace{-0.6em}\rule{\linewidth}{0.5pt}]
\titlespacing*{\section}{0pt}{5pt}{2pt}
\setlist[itemize]{leftmargin=1.2em,itemsep=0.2pt,topsep=1pt,parsep=0pt}
\newcommand{\role}[4]{\textbf{#1}\hfill\textbf{#2}\\\textit{#3}\hfill\textit{#4}}

\begin{document}\fontsize{9.8}{10.5}\selectfont

\begin{center}
{\LARGE\bfseries Amal Mehta}\\[3pt]
SF Bay Area \quad$\cdot$\quad \href{mailto:amal.mehta@gmail.com}{amal.mehta@gmail.com} \quad$\cdot$\quad \href{https://amalmehta.com/}{amalmehta.com}
% Saved for later, cut 2026-11-17, once ended the line above: \quad$\cdot$\quad \href{https://www.linkedin.com/in/amal-mehta-264413a9/}{LinkedIn} \quad$\cdot$\quad \href{https://github.com/amalmehta}{GitHub}
\end{center}

Research engineer in reinforcement learning and post-training, with four years in industry and deep learning research since 2020. I build each study end to end in PyTorch, from environment and training loop to held-out evals, seeds and a live demo.

\section*{Experience}

\role{Independent ML \& Robotics Research}{Sept 2025 -- Present}{Self-directed projects; details at \href{https://amalmehta.com/}{amalmehta.com}}{San Francisco Bay Area}
\begin{itemize}
\item \href{https://amalmehta.com/long_term_horizon_planning/}{\textbf{Reward hacking in RL-trained code agents:}} GRPO on visible tests raised pass rate 35\%\,$\to$\,58\% while true correctness fell 31\%\,$\to$\,2\%, as the agent deleted 98\% of untested branches. Ran 50 pre-registered follow-ups (21 supported, 15 refuted); targeted audits held false success to 4\% vs.\ 24\% for random audits at the same 10\% budget.
\item \href{https://github.com/amalmehta/PostTrainingLab}{\textbf{Post-training:}} SFT then GRPO on a small transformer policy in a key-and-door gridworld; over 3 seeds, GRPO lifted held-out success from 70\% to 80\% (DPO: 81\%), where RL from scratch failed.
\item \href{https://amalmehta.com/AlphaGoLite/}{\textbf{AlphaGo Lite:}} AlphaZero-style 9$\times$9 Go from scratch on one laptop CPU (107 generations, 17,120 self-play games); went from 0 to 32 wins in 40 games against GNU Go level 1. Runs in the browser (ONNX) and a Mac app (Core ML).
% Saved for later, cut 2026-11-24 so PhysicsAI stays the fullest section: \href{https://github.com/amalmehta/BCIProjects}{\textbf{Brain--computer interfaces:}} label-free decoding held R\textsuperscript{2} 0.64 on all 19 held-out days to day 138 (Wiener 0.35); 0.72 vs.\ 0.66 for a GRU without 2 remap days.
% Saved for later, cut 2026-11-19: \item \href{https://amalmehta.com/DroneAutonomy/}{\textbf{Meta-RL for a quadrotor:}} seven meta-learning methods compared on adapting to payload, wind and motor wear.
% Saved for later, cut 2026-11-19: \item \textbf{Tools:} \href{https://amalmehta.com/LLMUnderstander/}{LLM Understander}, an interactive \href{https://github.com/amalmehta/LLMUnderstander/releases/latest}{Mac app} on how LLMs are trained; \href{https://github.com/amalmehta/researcher_pal}{Researcher Pal}, Claude Code skills for research papers.
\end{itemize}

\role{Deep Learning Scientist}{Sept 2021 -- Aug 2025}{PhysicsAI (DARPA-funded)}{Pacifica, CA (remote)}
\begin{itemize}
\item Turned synthetic-data generation for drone perception into a closed RL loop: policies (TD3, PPO, DDPG) set an image generator's parameters, a spatiotemporal detector retrained on its images, and its accuracy on real drone imagery was the reward.
\item Built the PEARL meta-RL codebase (probabilistic task embeddings over TD3) for a fighter-jet simulation, so an agent adapts to a new task in a few tries; ran its multi-week GPU training.
% Saved for later, before 2026-11-24: Built the PEARL meta-RL codebase (probabilistic task embeddings over TD3) so the agent adapts to a new generator task in a few tries; ran its multi-week GPU training, and the method was also applied to a fighter-jet simulation.
\item Implemented gradient-reversal domain adaptation from scratch to hold detector accuracy with scarce real labels (10--15\,GB, mostly synthetic).
\item Built and benchmarked GAN synthesis pipelines (DCGAN, InfoGAN, FineGAN) that produced the detectors' training data.
\end{itemize}

\role{Deep RL Researcher}{Aug 2023 -- Mar 2024}{Hybrid Robotics Lab, UC Berkeley --- adv.\ Prof.\ Koushil Sreenath}{Berkeley, CA}
\begin{itemize}
\item Trained a single multi-stage PPO policy in IsaacGym for a Unitree Go1 quadruped to track, approach, and shoot a soccer ball at a target, replacing the multi-policy hierarchical framework used in prior work; $\sim$60\% kicking accuracy on hardware.
\item Designed the three-stage reward and CPG-based gait tracking; prototyped RLPD to fine-tune on real robot data.
\end{itemize}

\role{Research Intern}{May 2021 -- Aug 2021}{IBM Research}{Almaden, CA}
\begin{itemize}
\item Built an NLP pipeline (Watson NLP, human in the loop) that extracts terms from unstructured contracts and checks them against invoices to catch spend leakage, with a web app that reports the mismatches; patent submitted.
\end{itemize}

\role{Research Assistant}{June 2020 -- Dec 2022}{Video \& Image Processing Lab, BAIR --- adv.\ Prof.\ Avideh Zakhor}{UC Berkeley}
\begin{itemize}
\item Built a deep CNN segmentation pipeline for invasive melanoma in whole-slide histopathology (F1 0.44 vs.\ 0.14 prior best): patch-wise training and inference on gigapixel slides, detail-preserving downsampling, multi-GPU PyTorch/CUDA.
\item Published at SPIE Medical Imaging 2023 and USCAP 2022 (Stowell-Orbison Award Poster Session); mentored master's students extending it to vision transformers.
% Saved for later, the single bullet before 2026-11-19: Developed a deep CNN segmentation pipeline for invasive melanoma in large-scale histopathological images (F1 0.44 vs.\ 0.14 for the prior best); published at SPIE Medical Imaging 2023; Stowell-Orbison Award Poster Session at USCAP 2022.
\end{itemize}

\role{Researcher}{June 2020 -- Aug 2020}{Abbasi-Asl Lab, UCSF --- adv.\ Prof.\ Reza Abbasi-Asl}{San Francisco, CA}
\begin{itemize}
\item Built empirical mode decomposition (EMD) pipelines to align 2D and 3D brain scans across visits and to a common template.
\end{itemize}

\section*{Education}

\role{University of California, Berkeley}{2021}{B.A.\ Computer Science --- Regents Scholar}{Berkeley, CA}
% Saved for later, cut 2026-11-19: \item CS285 (Deep RL) project: iterative retraining of PPO agents against adversarial policies.

\role{Stanford University}{Aug 2021 -- May 2023}{Graduate coursework --- Autonomous Systems \& Artificial Intelligence}{}
\begin{itemize}
\item CS231N (Deep Learning for Vision), CS330 (Meta-Learning), AA271A/B (Autonomy I \& II), AA273 (State Estimation).
% Saved for later, cut 2026-11-19: \item CS330 project: few-shot imitation learning; implemented RL\textsuperscript{2} and inverse soft Q-learning in Meta-World.
\end{itemize}

\section*{Technical Skills}
\begin{tabular}{@{}p{0.9in}l}
\textbf{Languages} & Python, C++, Swift, JavaScript / TypeScript \\
\textbf{ML / RL} & PyTorch, TensorFlow, stable-baselines3, IsaacGym, OpenCV, CTC decoding \\
\textbf{Robotics} & ROS, Gazebo, Unitree Go1, CPG gait control, Kalman filtering, Shared autonomy \\
\textbf{Tools} & Docker, GitHub Actions CI, CUDA/multi-GPU, ONNX, Core ML, SwiftUI, Three.js \\
\end{tabular}

\section*{Publications}
\begin{enumerate}[leftmargin=1.4em,itemsep=1pt,topsep=1pt]
\item Shah, A., \textbf{Mehta, A.}, Wang, M., Neumann, N., McCalmont, T., Zakhor, A. ``Deep Learning Segmentation for Invasive Melanoma.'' \textit{SPIE Medical Imaging}, 2023.
\item Neumann, N., Wang, M., \textbf{Mehta, A.}, et al. ``Quantifying Invasive Melanoma Volume by Deep Learning Segmentation at the Pixel-Level.'' \textit{USCAP Annual Meeting}, 2022.
\item \textbf{Mehta, A.}, et al. ``Contract Information Extraction Using Watson NLP and Human-in-the-Loop.'' \textit{INFORMS Annual Meeting}, 2021.
\end{enumerate}

\end{document}
